Converts between supported unsigned integer and floating-point forms, or between the supported floating-point widths, according to the operand classes.

The operand classes select one of three conversion families.

\begin{itemize}
\item Fn to Fn \texttt{.S} converts D to S, while \texttt{.D} converts S to D.
\item Fn to Rn truncates toward zero and writes an unsigned 64-bit integer. A negative value or negative infinity
produces zero; positive overflow or positive infinity produces \texttt{0xffffffffffffffff}; NaN produces zero.
\item Rn to Fn converts the full unsigned 64-bit source range to the selected floating-point format using
\texttt{FSTATUS.RM}.
\end{itemize}
An invalid Fn-to-Rn conversion generates NV without OF or NX. A valid conversion that discards a fractional part
generates NX. Fn-to-Fn and Rn-to-Fn conversion use the common floating-point rounding, floating-point
exception-cause generation, condition-enable testing, and commit rules.

The operand classes select one of three conversion families.

\begin{itemize}
\item Fn to Fn \texttt{.S} converts D to S, while \texttt{.D} converts S to D.
\item Fn to Rn truncates toward zero and writes an unsigned 64-bit integer. A negative value or negative infinity
produces zero; positive overflow or positive infinity produces \texttt{0xffffffffffffffff}; NaN produces zero.
\item Rn to Fn converts the full unsigned 64-bit source range to the selected floating-point format using
\texttt{FSTATUS.RM}.
\end{itemize}
An invalid Fn-to-Rn conversion generates NV without OF or NX. A valid conversion that discards a fractional part
generates NX. Fn-to-Fn and Rn-to-Fn conversion use the common floating-point rounding, floating-point
exception-cause generation, condition-enable testing, and commit rules.
